% The lower bound at general margin: what is missing, what is not.

\section{Toward the matching lower bound}\label{sec:lower}

Krauth and M\'ezard conjectured the value of the limiting ratio
$M_N(\kappa)/N$ for every margin, not only at
$\kappa=0$~\cite{krauth1989storage}; the constant $0.833$ is the
zero-margin point of the conjectured curve $\alpha_\star(\kappa)$.
At $\kappa=0$ the conjecture is resolved: Ding and
Sun~\cite{ding2019capacity} proved $\alpha_\star$ is a lower bound
with positive probability, the sharp-threshold sequence of Nakajima
and Sun~\cite{nakajimasun2022} upgrades positive probability to
convergence in probability, Huang~\cite{huang2024capacity} proved the
matching upper bound, and the numerical hypotheses of both proofs are
certified~\cite{shmalo2026capacity}.  Away from $\kappa=0$ the upper
half is Theorem~\ref{thm:main}; the lower half has, at present, no
analytic theorem to certify.  This section delimits the gap: a lemma
on why unconditioned moment methods cannot close it at nonnegative
margins, an embedded bound that does give nonzero-margin capacity
lower bounds on the negative side, and a numerical verification that
the $\kappa$-dependent inputs of the Ding--Sun program hold at
five certified margins and at $\kappa=0$.

\subsection{Plain second moments fail at every margin}

Write $Z=Z_{N,\kappa}$ for the number of $w\in\{-1,+1\}^N$ satisfying
all $M=\lceil\alpha N\rceil$ constraints, and
$q_\rho(\kappa)=\Prob(X\ge\kappa,\,Y\ge\kappa)$ for $(X,Y)$ bivariate
standard Gaussian with correlation $\rho$.

\begin{lemma}\label{lem:secondmoment}
For every $\kappa\in\R$ and $\alpha>0$,
\[
\liminf_{N\to\infty}\frac1N\log\frac{\E[Z^2]}{(\E Z)^2}
\;\ge\; c(\alpha,\kappa)\;>\;0 .
\]
Consequently the guarantee of the Paley--Zygmund bound,
$\Prob(Z>0)\ge(\E Z)^2/\E[Z^2]$, is at most $e^{-cN+o(N)}$, and no
bound depending on $\E Z$ and
$\E[Z^2]$ alone can certify more, since a distribution placing mass
$(\E Z)^2/\E[Z^2]$ at $\E[Z^2]/\E Z$ and the rest at zero matches
both moments.  For $\kappa\ge0$ the same conclusion holds for every
weighted count $Z_f=\sum_w\prod_a f(\langle g_a,w\rangle/\sqrt N)$
with $f\ge0$ supported on $[\kappa,\infty)$ and
$0<\int f\,d\Phi<\infty$.
\end{lemma}

\begin{proof}
$\E Z=2^N\Psi(\kappa)^M$ by independence over constraints.  For a
pair $w,w'$ at overlap $\rho$ the constraint fields are bivariate
Gaussian with correlation $\rho$, so counting ordered pairs by
Hamming distance $k$,
\[
\E[Z^2]=2^N\sum_{k=0}^N\binom Nk\,q_{1-2k/N}(\kappa)^M .
\]
Fix $\rho_0\in(0,1)$, let $k^*=\lceil N(1-\rho_0)/2\rceil$ and
$\rho^*=1-2k^*/N$.  Keeping the single term $k^*$ and using
$\binom Nk\ge e^{N H(k/N)}/(N+1)$ with $H$ the natural-log binary
entropy ($H(p)=H(1-p)$, so the exponent may be written with
$(1+\rho_0)/2$),
\[
\frac1N\log\frac{\E[Z^2]}{(\E Z)^2}
\ge\underbrace{H\Bigl(\tfrac{1+\rho_0}2\Bigr)-\log2
+\alpha\bigl[\log q_{\rho_0}(\kappa)-2\log\Psi(\kappa)\bigr]}
_{=:\,\psi_{\alpha,\kappa}(\rho_0)}\;+\;O(N^{-1}\log N),
\]
the error, of order $N^{-1}\log N$, collecting the binomial
prefactor, the $|\rho^*-\rho_0|\le
2/N$ continuity of $\rho\mapsto\log q_\rho$ (differentiable, below),
and $|M-\alpha N|\le1$.  Now $q_0=\Psi(\kappa)^2$ and $H(1/2)=\log2$
give $\psi_{\alpha,\kappa}(0)=0$.  Reflecting the orthant to a joint
CDF, $q_\rho(\kappa)=\Phi_2(-\kappa,-\kappa;\rho)$, and Plackett's
identity $\partial\Phi_2/\partial\rho=\varphi_2$ with the symmetry
$\varphi_2(-\kappa,-\kappa;\rho)=\varphi_2(\kappa,\kappa;\rho)$ give
$\partial q_\rho/\partial\rho=\varphi_2(\kappa,\kappa;\rho)$ on
$(-1,1)$; at $\rho=0$ this is $\varphi(\kappa)^2$.  Since
$\frac{d}{d\rho}H(\tfrac{1+\rho}2)=\tfrac12\log\tfrac{1-\rho}{1+\rho}$
vanishes at $0$,
\[
\psi_{\alpha,\kappa}'(0)=\alpha\,
\frac{\varphi(\kappa)^2}{\Psi(\kappa)^2}\;>\;0 ,
\]
so $\psi_{\alpha,\kappa}(\rho_0)\ge\rho_0\,\psi'(0)/2>0$ for all
small $\rho_0>0$; fix such a $\rho_0$ and set
$c=\psi_{\alpha,\kappa}(\rho_0)$.  For the weighted counts, the diagonal
term of $\E[Z_f^2]$ is $2^N\E[f(X)^2]^M$, so either
$\E[Z_f^2]=\infty$ and the claim is trivial, or $f\in L^2(\Phi)$ and
Mehler's expansion gives
$\E[f(X)f(Y)]=\sum_{j\ge0}\rho^j\langle f,\mathrm{He}_j\rangle^2/j!$,
whose $\rho$-derivative at $0$ is $(\int xf\varphi\,dx)^2$; for
$f\ge0$ supported on $[\kappa,\infty)$ with $\kappa\ge0$ the
integrand $xf(x)\varphi(x)$ is nonnegative, and it vanishes almost
everywhere only if $f$ vanishes almost everywhere off the origin,
which $0<\int f\,d\Phi$ excludes; the argument
then runs unchanged.
\end{proof}

The mechanism is the asymmetry: pairs of solutions at small positive
overlap are exponentially favored over independent pairs at every
margin, which is what the AMP-conditioned second moment
of~\cite{ding2019capacity} is built to circumvent at $\kappa=0$, and
what any lower-bound proof must circumvent at every other margin.
The zero-margin instance of this failure is classical
(see the introduction of~\cite{ding2019capacity}); the content here
is the uniform mechanism and the closure over nonnegative-margin
weights.  For $\kappa<0$ the weighted closure genuinely fails, and
profitably so: a window weight symmetric about the origin has zero
linear Mehler coefficient, which is the next observation.

\subsection{A capacity lower bound at negative margins}

\begin{prop}\label{prop:embedded}
Let $\kappa<0$ and let $\alpha^{\mathrm{sym}}_c(K)=
\log2/(-\log\Prob(|X|\le K))$ denote the symmetric binary perceptron
threshold at width $K$.  For every
$\alpha<\alpha^{\mathrm{sym}}_c(|\kappa|)$, with probability
$1-o(1)$ the margin-$\kappa$ perceptron admits a solution:
$M_N(\kappa)/N\ge\alpha$ whp.
\end{prop}

\begin{proof}
If $|\langle g_a,w\rangle|/\sqrt N\le-\kappa$ for all $a$ then in
particular $\langle g_a,w\rangle/\sqrt N\ge\kappa$, so the symmetric
solution set at width $|\kappa|$ embeds in the margin-$\kappa$
solution set.  The symmetric threshold is
$\alpha^{\mathrm{sym}}_c$~\cite{aubin2019storage}; for Gaussian
disorder the satisfiability side below it holds with high
probability by Altschuler's critical-window
theorem~\cite{altschuler2022}.  (The lognormal limit of Abbe, Li,
and Sly~\cite{abbelisly2021} is proved for Rademacher disorder, and
Perkins and Xu~\cite{perkinsxu2021} treat the Gaussian case under a
single-critical-point hypothesis, so neither alone covers the model
here.)
\end{proof}

The bound is far from sharp --- at $\kappa=-0.05$ it gives
$0.215$ against the predicted $0.889$, at $\kappa=-0.45$ it gives
$0.655$ against $1.558$ --- and it is not alone: Li, Schramm, and
Zhou~\cite{lsz2024} prove lower bounds on the satisfiability
threshold of the same one-sided perceptron at negative margin by a
different route, with the sharpest statements in the
$\kappa\to-\infty$ regime.  What the embedding adds is its
elementary proof and margin-explicit values at the certified
points, strictly positive on the whole negative axis.

\subsection{The Ding--Sun condition at general margin}

The construction underlying the $\kappa=0$ lower bound
of~\cite{ding2019capacity} is written at general margin throughout
their Section~2: their equations (34)--(48) define, for parameters
$(\kappa_{\mathrm{all}},\alpha_{\mathrm{all}})$, an explicit
univariate function $S_{\kappa,\alpha}(\lambda)$ on
$[\lambda_{\min},1]$ built from the fixed point $(q_\star,\psi_\star)$
--- an entropy term through the two-spin measure $P_{H,D}$, a
constraint term through the conditional Mills law
$\varphi_\xi$, and a convex $\inf_s$ correction --- and their
Condition~1.2 is the statement $S_{0,\alpha_\star}(\lambda)<0$ for
$\lambda\notin\{0,1\}$.  The $\kappa$-analog of that hypothesis is
therefore already defined in their own notation, and every
$\kappa$-dependent ingredient it consumes is an object this paper
certifies: the fixed point, its uniqueness rectangle, and the
crossing (Theorem~\ref{thm:strip} and the deep enclosures of
Theorem~\ref{thm:main}).

We evaluated $S_{\kappa,\alpha_\star(\kappa)}$ in floating point
(the driver \texttt{ds\_condition.py} uses Gauss--Hermite and
conditional Gauss--Legendre quadrature) at five of the certified
margins and at $\kappa=0$, checking the six endpoint identities
$H(0)=P(0)=A(0)=0$, $H(1)=-H_\star$, $P(1)=-P_\star$,
and $G_\star=H_\star+P_\star\approx0$ (with $H_\star$, $P_\star$ the
entropy and constraint halves of the free energy), and that the
conditional-mean check resolves at $10^{-12}$ or better at every
margin, as implementation controls.  The result is uniform
(Figure~\ref{fig:dscondition}): at every margin the curve is
strictly negative off $\{0,1\}$, at most
$-5\cdot10^{-5}$ outside windows of half-width $0.02$ around the two
roots (inside them it climbs to its maxima at the roots, reaching
$-1.6\cdot10^{-6}$ at the closest grid points), the maximum at
$\lambda=0$ is degenerate exactly as
at $\kappa=0$, and the whole family collapses nearly onto a single
shape, with $\lambda_{\min}$ drifting from $-0.441$ at
$\kappa=-0.45$ to $-0.421$ at $\kappa=0.13$.  These are diagnostics,
not certificates; certifying $S_{\kappa,\alpha_\star}<0$ off the
roots is a univariate problem of exactly the species the
Condition~\ref{cond:2var} machinery already handles, and we expect it
to be far easier than the two-variable certification carried out
here.

\begin{figure}[t]
\centering
\includegraphics[width=0.85\linewidth]{ds_condition.png}
\caption{The general-margin Ding--Sun lower-bound condition,
evaluated at five of the certified margins and $\kappa=0$: the
univariate function $S_{\kappa,\alpha_\star(\kappa)}$ of their
equation (48), whose nonpositivity off $\{0,1\}$ is the numerical
hypothesis of their second-moment argument.  All six curves are
negative off the two roots and nearly coincide.}
\label{fig:dscondition}
\end{figure}

What separates these verified inputs from a general-margin lower
bound is the probabilistic body of~\cite{ding2019capacity}: the
AMP-conditioned moment computation of their Sections~2--7, whose
displays carry $\kappa_{\mathrm{all}}$ explicitly but whose proofs
are executed at $\kappa=0$.  An audit of that body would need, beyond
the certified inputs above: a priori bounds at margin (the upper one
is free at every $\kappa$ from the first moment, since
$\E Z\to0$ exponentially above the annealed threshold
$-\log2/\log\Psi(\kappa)$; the lower one is
Proposition~\ref{prop:embedded} for $\kappa<0$ and open for
$\kappa>0$), the state-evolution inputs (generic in the
nonlinearity, hence in $\kappa$), and the sharp-threshold transfer
of~\cite{nakajimasun2022} if convergence in probability rather than
positive probability is wanted.  We record this as the concrete
route to the exact threshold at the certified margins: the numerical
side is done at six points and the analytic side is a
$\kappa$-generalization audit of an existing proof, not a new
construction.

That audit is now carried out for their Sections~2--8.  Sections
2--6 generalize verbatim: the margin rides through the TAP
conditioning, the change of measure, and the admissibility and
volume estimates, entering only through $(q_\star,\psi_\star)$ and
the shifted constraint kernels.  In Section~7 one proof step is
genuinely $\kappa=0$: the monotonicity argument of their Lemma~7.3
regroups the Gaussian under a reflection available only at zero
margin; its conclusion is exactly the $\partial R/\partial q$
positivity that Lemma~\ref{lem:dR} proves at every $\kappa$, so the
repaired statement is what our certificates consume.  In Section~8
the analytic skeletons --- the sandwich bounds on the entropy
parametrization, the square-root singularity of the constraint term
at $\lambda=1$, and the special-function inequalities behind both
--- hold at every margin, while the five numerical constants they
instantiate at $\kappa=0$ become per-margin certified integrals over
the located $(q_\star,\psi_\star)$ boxes.  The grid bounds of their
Section~9 are not ported at all: at the certified margins their role
is played by the certificates of this paper.  Of the endpoint
package their assembled proof consumes, the derivative divergence at
$\lambda=1$ and the second-derivative sign at $\lambda=0$ become
per-margin certified objects once the ported constants of their
Section~8 land (their status is tracked in the audit file); the
value at $\lambda=1$ is analytic.  Their Section~2 already states
$P(1)=-P_\star$ at general margin; what the following lemma adds is
the full tilted family and the assembly --- $I_s\to0$ for every
$s$, hence $A(1)=0$ and the sum --- with the dominated-convergence
argument written out.

\begin{lemma}\label{lem:S1}
For every $\kappa$ and every $\alpha$ in a certified window,
$S_{\kappa,\alpha}(1)=-G_\star$, where
$G_\star=H_\star+P_\star$ is the first-moment exponent at
$(\kappa,\alpha)$.  Moreover the $s$-tilted constraint term
satisfies $I_s(\lambda)\to0$ as $\lambda\uparrow1$ for every fixed
$s$.  Here and throughout, values at $\lambda=1$ denote the
$\lambda\uparrow1$ limits: the integral formulas are singular at
$\lambda=1$ itself, and this is the convention under which
Ding--Sun consume the endpoint.
\end{lemma}

\begin{proof}
The explicit middle term of the constraint half carries the factor
$(1-\lambda)$ and vanishes at $\lambda=1$.  For the integral term,
condition on the first copy's variables: the second copy's field
$\nu$ follows the Mills law $\varphi_\xi$ on the ray $\nu>\xi$,
whose support does not depend on $\lambda$, and the constraint
factor is $\log\Psi$ of
\[
\frac{\xi-\lambda\nu}{\sqrt{1-\lambda^2}}
-\frac{E(\xi)\,s}{\psi_\star^{1/2}\sqrt{1-q_\star}}
=\frac{\xi(1-\lambda)-\lambda(\nu-\xi)}{\sqrt{1-\lambda^2}}
-\frac{E(\xi)\,s}{\psi_\star^{1/2}\sqrt{1-q_\star}}\,.
\]
On the open support $\nu-\xi>0$ the first fraction tends to
$-\infty$ as $\lambda\uparrow1$, and a fixed finite $s$-shift does
not affect the divergence, so $\log\Psi\to0$ pointwise for
$\varphi_\xi$-almost every $\nu$.  Domination needs no case split:
for every $0\le\lambda<1$ and every $\nu\ge\xi$ the argument is at
most
\[
\xi_+\sqrt{\frac{1-\lambda}{1+\lambda}}
+\frac{E(\xi)\,|s|}{\psi_\star^{1/2}\sqrt{1-q_\star}}
\;\le\; |\xi| + \frac{(|\xi|+1)\,|s|}{\psi_\star^{1/2}
\sqrt{1-q_\star}} \;=:\; M(z)\,,
\]
using $E(\xi)\le|\xi|+1$ (Lemma~\ref{lem:mills} gives
$E'\in(0,1)$, so $E$ is increasing with unit-bounded slope, and
$E(0)=\sqrt{2/\pi}<1$); since $|\log\Psi|$ is increasing and
$\xi=(\kappa-q_\star^{1/2}z)/\sqrt{1-q_\star}$ is affine in $z$,
so that $M(z)$ is affine in $|z|$, the
integrand is bounded by $|\log\Psi(M(z))|\le C_s(1+z^2)$ for every
$\lambda$ and $\nu$ at once, and this majorant is
$\varphi(z)\,dz$-integrable because it does not depend on $\nu$.
Dominated convergence gives $I_s(\lambda)\to0$ for every
$s$; in particular $P(1)=-P_\star$.

For the $\inf_s$ correction, write $B_\lambda(s)$ for the bracket,
so $A(\lambda)=\inf_s B_\lambda(s)$, and note that fixed-$s$
convergence alone does not pass through an infimum: the limit and
the infimum must be exchanged with care.  Three facts do it.
First, $B_\lambda(0)=0$ identically ($I_0=I$), so $A(\lambda)\le0$.
Second, $s\mapsto B_\lambda(s)$ is convex --- the convexity of the
$s$-tilt recorded with their Lemma~5.7.  Third, the dominated
convergence above runs uniformly over $|s|\le1$: the majorant at
$s=1$ dominates every $|s|\le1$ since $M(z)$ is increasing in
$|s|$, and $\sup_{|s|\le1}$ of the integrand still vanishes
pointwise because $|\log\Psi|$ is increasing and the supremum of
the argument over $|s|\le1$ also tends to $-\infty$; hence
$\sup_{|s|\le1}|I_s(\lambda)-I(\lambda)|\to0$.  Now fix $\lambda$
close enough to $1$ that $B_\lambda(\pm1)\ge\tfrac14$, possible by
the fixed-$s$ limits $B_\lambda(\pm1)\to\tfrac12$.  For $|s|>1$,
convexity through the points $0$ and $s$ gives
$B_\lambda(s/|s|)\le B_\lambda(s)/|s|$, so $B_\lambda(s)<0$ would
force $B_\lambda(\pm1)<0$; hence every $s$ with $B_\lambda(s)<0$
lies in $[-1,1]$, and since $B_\lambda(0)=0$,
$A(\lambda)=\inf_{|s|\le1}B_\lambda(s)$.  On that interval,
$B_\lambda(s)\ge -c_\lambda-\sup_{|s|\le1}|I_s(\lambda)-I(\lambda)|$
with $c_\lambda=\psi_\star^{1/2}\sqrt{1-q_\star}\,
\sqrt{(1-\lambda)/(1+\lambda)}\to0$ the linear coefficient, and
together with $A(\lambda)\le0$ this gives $A(\lambda)\to0$, that
is, $A(1)=0$.  The entropy half is prior-side, with
$\kappa$ entering only through $(q_\star,\psi_\star)$, and
$H(1)=-H_\star$ as at zero margin.  Summing,
$S(1)=-P_\star-H_\star=-G_\star$.
\end{proof}

For the second-moment argument only $S(1)\le-c<0$ is consumed, and
below the crossing $G_\star>0$ strictly, so the lemma supplies the
endpoint with margin.  What remains open on this route is
quantitative rather than structural: the per-margin constants of
Section~8 and the slack they leave against our wider
$(q_\star,\psi_\star)$ enclosures, both objects of the same
certified species as everything above.
